\DocSection{Expected-Return Policy Target}
\label{sec:target}

For inverse temperature \(\beta>0\), define the target measure
\begin{equation}
  \Target(d\policy)
  = \frac{\exp\!\left(\beta\ExpectedReturn_H(\policy)\right)}
         {Z_{\beta}}\,\mu(d\policy),
  \qquad
  Z_{\beta}
  = \int_{\PolicySpace}
      \exp\!\left(\beta\ExpectedReturn_H(\policy)\right)\mu(d\policy).
  \label{eq:policy-target}
\end{equation}
Bounded rewards and a finite reference measure make \(Z_\beta\) finite.
The parameter \(\beta\) carries the inverse units of reward, so a change of
reward units must be accompanied by the inverse change in \(\beta\).

\begin{proposition}[Mode preservation]
Every mode of \(\Target\) with respect to a uniform finite reference measure
maximizes \(\ExpectedReturn_H\), and every expected-return maximizer is a mode.
\end{proposition}

\begin{proof}
The exponential function is strictly increasing, and the normalizing constant
does not depend on \(\policy\).
\end{proof}

The order of expectation and exponentiation is essential. In general,
\begin{equation}
  \exp\!\left(\beta\E[\Return_H]\right)
  \neq
  \E\!\left[\exp(\beta\Return_H)\right].
  \label{eq:order-matters}
\end{equation}
The right-hand side defines an entropic-risk quantity and places additional
weight on return variability. It is not an estimator identity for the target
in Equation~\eqref{eq:policy-target}.

The target also has the variational characterization
\begin{equation}
  \log Z_\beta
  = \sup_Q
    \left\{
      \beta\E_{\policy\sim Q}[\ExpectedReturn_H(\policy)]
      -\KL(Q\|\mu)
    \right\}.
  \label{eq:gibbs-variational}
\end{equation}
Thus the regularization acts on a distribution over complete policies. It is
not the same as adding independent action entropy at every visited state.

